% Français 9117, batch 1 — Gallica views 1–20.
% Pass: Opus 5 (claude-opus-5), 2026-09-14 — first pass, unchecked against
% the views by a human.
%
% What the twenty views hold. The opening of the volume, then the opening of
% the work: the working manuscript of the « Considérations générales sur
% l'état des sciences et des lettres », in French prose, no mathematics set
% as formulae. It is a draft for the press rather than a fair copy: a regular
% hand, but struck, overwritten and rewritten between the lines, some
% paragraphs opened by a square bracket, and three names in another, smaller
% hand in the margins or at the head (« Gutmann » v. 9, « Cherenoz » v. 17,
% « Deraihe » v. 19, the last two uncertain) — very likely the compositors'
% names marking where each took up the copy; that is this pass's inference.
%
%   v. 1       Microfilm target (« Ce document a été microfilmé tel qu'il a
%              été relié »). Not transcribed.
%   v. 2–4     Marbled boards with the label « FR. 9117 », the inside of the
%              binding, a blank flyleaf. Not transcribed.
%   v. 5       The library's title leaf (folio « 1 »): a note in a
%              librarian's hand — shelfmark « Supplt fr. 2209 », « Manuscrit
%              autographe de Mademoiselle Germain », the title of the 1833
%              edition, « donné par M. Lherbette député & neveu de l'auteur »,
%              an accession reference « R.B. N° 106 4. 1833 ». Not transcribed;
%              it is used only in a note to identify the signature of v. 7.
%   v. 6       Blank.
%   v. 7       Folio 2, a small leaf on a stub: the covering letter, not in
%              Germain's hand, with which Lherbette sends the manuscript and a
%              copy of the posthumous book to Champollion at the Bibliothèque
%              royale, dated only « 7bre. ». Transcribed, as the
%              correspondence volumes transcribe letters to and about her.
%   v. 8       Its verso: the address « Monsieur Champollion Conservateur à la
%              Bibliothèque Royale » (given in a note under v. 7), and blank
%              paper.
%   v. 9–10    « Ch. I. » and its heading « Comment les sciences et les
%              lettres sont dominées par un sentiment qui leur est commun »:
%              the likeness of all the works of the mind, the laws of order
%              and proportion, whether the type reveals the conditions of
%              being or only how we understand the possible; poem, natural
%              history, exact sciences. V. 10 breaks off mid-sentence and
%              v. 13 finishes it.
%   v. 11      A smaller leaf pasted in (folio 5), marked « + »: an addition —
%              a stroke of genius pleases in every field for one reason, it
%              discloses unforeseen relations. The same cross stands at the
%              foot of v. 13, which is where it is meant to go; it is left in
%              the order of the views.
%   v. 12      Its verso, blank (show-through only).
%   v. 13      The exact sciences as the closest realisation of the type;
%              imagination and demonstration seem apart, yet one method guides
%              both. The last sentence is heavily corrected and not settled.
%   v. 14      Blank verso.
%   v. 15–18   The comparison of the poet and the geometer at each stage of
%              composition: first essays, the one simple idea chosen, the plan
%              and unity of the work, filling the frame, then turning judges
%              of their own work — episodes and notes cut out, style,
%              corrections of taste. V. 16 opens with a paragraph crossed out
%              whole and rewritten below it.
%   v. 19      The corrections proper to « la langue des calculs »: the choice
%              of characters, formulas in place of the phrase, analysis
%              speaking to the eyes, tact as taste applied to objects thought
%              foreign to it; then the summing-up of the chapter so far.
%   v. 20      Why the operations of the mind, one and the same, have
%              received different names according to their subjects; the poet
%              will not tell of the discussions behind his emblems, nor the
%              man of genius of his wanderings. It breaks off mid-sentence,
%              « qui devoit le conduire à la », at the foot of the page.
%
% The numbering on the leaves. Two series, both read on the views.
% (1) Top right of the rectos, a thin, light stroke unlike the text's ink —
% taken to be the library's foliation, though a grayscale scan cannot show
% pencil for certain: « 1 » v. 5, « 2 » v. 7, « 4 » v. 9, « 5 » v. 11,
% « 6 » v. 13, « 7 » v. 15, « 8 » v. 17, « 9 » v. 19. The « 4 » of v. 9 is a
% looped figure, the same form as the 4 of « 40 » at v. 81; leaf 3, which
% would be the address panel at v. 8, carries no visible number. \folio{N}
% is written bare, on rectos only, never inferred. (2) Top left of every written page, in ink, the author's (or the
% copy's) page series: « 1 » v. 9, « 2 » v. 10, « 3 » v. 13, « 5 » v. 15,
% « 6 » v. 16, « 7 » v. 17, a figure under a blot at v. 18 (presumably 8),
% « 9 » v. 19, « 10 » v. 20. No « 4 » was found; the pasted leaf v. 11 bears
% only « + ». These are recorded in the notes, not as \folio{}.
%
% Notation. None: there is no formula in the batch. Spelling of the time is
% kept (« connoître », « étoit », « longtems », « différens », « sentimens »,
% « poëte », « apperçoive », « assujetie » at v. 9 and « assujetti » later);
% the ampersand is set \&. Her deletions are \struck{}; words written
% between the lines in replacement follow the struck words they replace.
% Where corrections pile up in layers (v. 16, v. 17, v. 19, v. 20) the
% uncancelled reading is given and the notes say what could not be ordered.
%
% What could not be read (\ill{}): v. 10 the struck first word; v. 13 one
% word (« la littérature la plus \ill{} »); v. 15 one struck word; v. 16
% seven — four in the crossed-out paragraph, three in struck interlinear
% corrections; v. 19 two, inside a layered deletion; v. 20 two, one
% overwritten word in the interlinear correction and one struck word before
% « davantage ». Uncertain readings are flagged \uncertain{} in place.
%
% Nothing here was checked against any published transcription.
\documentclass[11pt,a4paper]{article}
% The preamble shared by every transcription of the mathematicians' manuscripts in Gallica.
%
% It rests on one idea: **uncertainty compounds, it does not resolve.** A
% transcription that smooths over a doubtful word has destroyed the only thing
% it was made for — a reader must be able to tell, at every point, what was on
% the page and what was supplied. The apparatus macros below are therefore the
% critical apparatus, not decoration.
%
% The same commands are recognised by `scripts/render.mjs`, which produces the
% site's reading view, and by `scripts/tei.mjs`. A macro added here must be
% added there too, or rendering fails loudly — which is the wanted behaviour.
%
% Comments here are English, like the rest of the repository. The documents
% this preamble sets are French, and that is deliberate: see the skills.

\ProvidesPackage{germain}[2026/09/15 Transcriptions of mathematicians' manuscripts in Gallica]

\usepackage{amsmath,amssymb,amsthm}
\usepackage{mathrsfs}
% Kept from the parent preamble so the renderer's diagram support stays
% available; these manuscripts draw no commutative diagrams, and nothing here uses it.
\usepackage{tikz-cd}
\usepackage[english,french]{babel}
\usepackage{fontspec}

% --- Greek ------------------------------------------------------------------
% The text face stays fontspec's default, Latin Modern, which has no Greek:
% XeTeX drops every glyph it lacks with a line in the log and nothing on the
% page, and the Greek words the manuscripts quote vanished from the PDFs. So
% the Greek and Greek Extended blocks get a XeTeX character class of their own,
% and entering or leaving a run of it switches the family to CMU Serif — the
% same Computer Modern design, with polytonic Greek — and back. Only the
% family changes: a Greek word in \textit or \textbf keeps its shape and series.
% The transitions to and from every other class carry the tokens that class
% has towards an ordinary letter, so French spacing before « ; » or inside
% « » still holds after a Greek word. No grouping, so no brace or \endgroup
% can come out unbalanced; maths is left alone, since XeTeX inserts nothing
% there. Kept identical in `exercices.sty`.
\makeatletter
\newfontfamily\germainGreekfont{cmun}[
  Extension=.otf, NFSSFamily=germaingreek,
  UprightFont=*rm, ItalicFont=*ti, SlantedFont=*sl,
  BoldFont=*bx, BoldItalicFont=*bi, BoldSlantedFont=*bl]
\newXeTeXintercharclass\germain@greek
\def\germain@greekrange#1#2{%
  \@tempcnta=#1\relax
  \loop
    \XeTeXcharclass\@tempcnta=\germain@greek
    \ifnum\@tempcnta<#2\advance\@tempcnta\@ne\repeat}
\germain@greekrange{"0370}{"03FF}
\germain@greekrange{"1F00}{"1FFF}
\def\germain@greekprev{\rmdefault}
\def\germain@greekin{%
  \edef\germain@greekprev{\f@family}\fontfamily{germaingreek}\selectfont}
\def\germain@greekout{\fontfamily{\germain@greekprev}\selectfont}
\def\germain@greekjoin{%
  \XeTeXinterchartokenstate=\@ne
  \@tempcnta=\z@
  \loop
    \ifnum\@tempcnta=\germain@greek\else
      \XeTeXinterchartoks\germain@greek\@tempcnta=\expandafter{%
        \expandafter\germain@greekout\the\XeTeXinterchartoks\z@\@tempcnta}%
      \XeTeXinterchartoks\@tempcnta\germain@greek=\expandafter{%
        \the\XeTeXinterchartoks\@tempcnta\z@\germain@greekin}%
    \fi
    \ifnum\@tempcnta<4095\advance\@tempcnta\@ne\repeat}
% babel-french sets its own classes, and saves and restores the interchar
% state, each time a language is selected: join again after it does.
\addto\extrasfrench{\germain@greekjoin}
\addto\extrasenglish{\germain@greekjoin}
\AtBeginDocument{\germain@greekjoin}
\makeatother

\usepackage{geometry}
\geometry{a4paper,margin=25mm,marginparwidth=28mm}
\usepackage{marginnote}
\usepackage{xcolor}
\usepackage[normalem]{ulem}
\setlength{\emergencystretch}{3em}
\usepackage{hyperref}
\hypersetup{colorlinks=true,linkcolor=black,urlcolor=black,pdfborder={0 0 0}}

% --- Batch metadata ---------------------------------------------------------
% Taken as they stand by the rendering script, which shows them at the head of
% the reading view. They fix what the file claims to cover. The macro names are
% the parent project's, so that the scripts stay shared; \folder here means the
% volume's slug — `fr-9115` — and \pages counts Gallica views.

\newcommand{\folder}[1]{\gdef\grothFolder{#1}}
\newcommand{\batch}[1]{\gdef\grothBatch{#1}}
\newcommand{\pages}[2]{\gdef\grothFirst{#1}\gdef\grothLast{#2}}
\newcommand{\foldertitle}[1]{\gdef\grothFoldertitle{#1}}

% The holder's own shelfmark, printed as they write it: « Français 9115 ».
\newcommand{\shelfmark}[1]{\gdef\grothShelfmark{#1}}

% The dating, copied verbatim from the holder's catalogue — never inferred
% here. « XIXe siècle » is the BnF's; a pass that dates a leaf from its content
% says so in a \note{}, as its own inference.
\newcommand{\dating}[1]{\gdef\grothDating{#1}}

% The status notice, printed in the title block and repeated as a diagonal
% watermark on every page of the PDF. The manuscripts are in the public
% domain; what this line declares is not a right but a quality — an
% unverified first pass by a machine. The skills write the line; a file
% without it gets no watermark, so leaving it out is visible in the output.
\newcommand{\watermark}[1]{\gdef\grothWatermark{#1}}

\gdef\grothFolder{?}\gdef\grothBatch{}\gdef\grothFirst{?}\gdef\grothLast{?}%
\gdef\grothFoldertitle{}\gdef\grothShelfmark{}\gdef\grothDating{}\gdef\grothWatermark{}

% --- The résumé -------------------------------------------------------------
\newenvironment{resume}
  {\small\setlength{\parskip}{0.45em}}
  {\par\medskip\hrule\medskip}

% --- Keywords ---------------------------------------------------------------
% In English, closing the modernised reading's résumé: the modern vocabulary
% under which to search for what the leaves construct. The manifest extracts
% them to tag the volume — this line is the single source of the tags.
\newcommand{\keywords}[1]{\par\smallskip\noindent
  {\footnotesize\textsc{Keywords} --- \textit{#1}}\par}

% --- The critical apparatus -------------------------------------------------

% The Gallica view number, in the margin. This is the anchor that lets a
% disagreement over a formula be settled by looking at *one* image rather than
% a volume of 750. The site uses it to turn the facsimile as the transcription
% is scrolled. A view is one scanned image: usually one side of a leaf.
\newcommand{\page}[1]{%
  \marginnote{\footnotesize\color{gray!70}\textsf{v.\,#1}}%
  \label{page:#1}\ignorespaces}

% The BnF's pencilled foliation, where the leaf carries it and a pass can read
% it: « 198r ». This is what the literature cites — « ff. 198r–208v » — and
% the one line that turns a citation into a view. Recorded, never inferred:
% a folio a pass cannot read is not written.
\newcommand{\folio}[1]{%
  \marginnote{\footnotesize\color{gray!50}\textsf{f.\,#1}}\ignorespaces}

% The modernised reading's coarser counterpart to \page{}: which views a
% section is drawn from.
\newcommand{\pagerange}[2]{%
  \marginnote{\footnotesize\color{gray!70}\textsf{v.\,#1--#2}}%
  \ignorespaces}

% Illegible: never guessed.
\newcommand{\ill}{\textcolor{red!60!black}{[\,\ldots\,]}}

% A doubtful reading: the word is given, and flagged as uncertain.
\newcommand{\uncertain}[1]{\uline{#1}}

% An editorial addition — a missing letter, an implied word.
\newcommand{\add}[1]{\textcolor{blue!50!black}{[#1]}}

% Struck out by the writer. What was crossed out is part of the document.
\newcommand{\struck}[1]{\sout{#1}}

% The transcriber's note, on the state of the page or the mathematics.
\newcommand{\note}[1]{\par\smallskip{\small\color{gray!60!black}\textit{[#1]}}\par\smallskip}

% A marginal note of hers, distinct from the body of the text.
\newcommand{\marginal}[1]{\par\smallskip\noindent\hspace{1em}%
  {\small\color{gray!60!black}#1}\par\smallskip}

% --- Title ------------------------------------------------------------------
% The title block is French, like the documents themselves.

\AddToHook{shipout/background}{%
  \ifx\grothWatermark\empty\else
    \begin{tikzpicture}[remember picture, overlay]
      \node[rotate=52, scale=3.4, align=center, text=gray!18,
            font=\sffamily\bfseries]
        at (current page.center) {\grothWatermark};
    \end{tikzpicture}%
  \fi}

\AtBeginDocument{%
  \begin{center}
    {\large\bfseries Manuscrits de mathématiciens (Gallica) ---
      \ifx\grothShelfmark\empty\grothFolder\else\grothShelfmark\fi}\\[2pt]
    {\itshape\grothFoldertitle}\\[4pt]
    {\small vues \grothFirst--\grothLast
      \ifx\grothBatch\empty\else\ (lot \grothBatch)\fi}\\[2pt]
    \ifx\grothDating\empty\else
      {\small\color{gray!60!black}Datation du catalogue : \grothDating}\\[2pt]
    \fi
    \ifx\grothWatermark\empty\else
      {\small\bfseries\color{red!45!gray}\grothWatermark}\\[2pt]
    \fi
    {\footnotesize\color{gray!60!black}Transcription établie d'après les images IIIF de Gallica.
     Source gallica.bnf.fr / Bibliothèque nationale de France.\\
     Les lectures incertaines sont soulignées, les illisibles notées [\,\ldots\,],
     les ajouts entre crochets ; \textsf{v.} numérote les vues de Gallica,
     \textsf{f.} la foliotation au crayon de la BnF.}
  \end{center}
  \vspace{1em}}

\endinput


\folder{fr-9117}
\shelfmark{Français 9117}
\batch{1}
\pages{1}{20}
\dating{1801-1900}
\watermark{Lecture automatique\\non vérifiée}
\foldertitle{« Considérations générales sur l'état des sciences et des lettres aux différentes époques de leur culture, » par Sophie GERMAIN.}

\begin{document}

\section{Lettre d'envoi du manuscrit à la Bibliothèque royale}

\page{7}\folio{2r}
\note{feuillet de plus petit format, monté sur onglet ; la main n'est pas celle de Sophie Germain. Timbre « Bibliothèque royale » non transcrit. L'adresse, au verso (vue 8) : « Monsieur Champollion Conservateur à la Bibliothèque Royale ».}

Monsieur,

D'après le désir que vous avez bien voulu m'en témoigner, j'ai l'honneur de
vous adresser, avec un exemplaire d'un ouvrage posthume de Mlle Germain, le
manuscrit que vous joindrez, si vous le jugez à propos, à ceux qui sont déjà
déposés à la Bibliothèque.

Agréez, je vous prie, l'assurance des sentimens distingués avec lesquels j'ai
l'honneur d'être,

Monsieur,

7bre.

Votre très humble \& obéissant serviteur

Lherbette

\note{signature lue d'après la note du feuillet de titre (vue 5), qui nomme « M. Lherbette député, \& neveu de l'auteur » comme donateur ; « 7bre. » est écrit à gauche de la formule finale, sans jour ni année.}

P.S. Me permettez-vous, Monsieur, de joindre ici un exemplaire pour
M. Letronne, qui, ne connaissant pas \uncertain{personnellement} Mlle Germain,
ne comprendrait pas l'envoi que je lui \uncertain{en} ferais directement\,?

\section{Ch. I.}

\page{9}\folio{4r}
\note{en tête du feuillet, à gauche, le chiffre « 1 » à l'encre ; au centre le titre « Ch. I. », souligné.}

Comment les sciences et les lettres sont dominées par un sentiment qui leur est
commun.

\note{dans la marge de gauche, au début du paragraphe suivant, le nom « Gutmann », souligné, d'une écriture plus petite ; le paragraphe est ouvert par un crochet.}

Lorsqu'on envisage sous un point de vue général les divers travaux de l'esprit
humain, on est frappé de leur similitude. Partout de certaines lois ont été
observées ; ou, si elles ne l'ont pas été, leur défaut s'est fait sentir, et
alors, soit que l'ouvrage renferme un corps de doctrine, soit qu'il ait été
destiné au simple amusement du lecteur, l'auteur n'a pas rempli les conditions
de la durée ; à la première curiosité, bientôt épuisée, succèdera un entier
oubli.

\note{« et » est ajouté dans l'interligne au-dessus d'un mot recouvert devant « alors ».}

Les lois dont nous parlons ont régi la pensée de l'homme longtems avant qu'il
ait eu le loisir de réfléchir. Le spectacle de l'univers en étoit empreint ; la
mémoire les a reproduites ; l'imagination, jusque dans ses caprices, leur est
demeurée \struck{soumise} assujetie ; plus tard elles ont servi de guide à la
raison.

S'il nous étoit donné de pénétrer la nature des choses, si les observations,
les réflexions, les théories qui composent notre richesse intellectuelle
n'étoient pas de l'homme, nous choisirions avec certitude entre ces deux
propositions : ou le type que nous trouvons en nous mêmes et dans les objets
extérieurs nous révèle les conditions de l'être ; ou ce type, nous appartenant
en propre, atteste seulement la manière dont nous pouvons comprendre les
possibles.

Cette haute connoissance nous est à jamais interdite ; mais en nous bornant à
chercher comment un sentiment profond d'ordre et de proportions devient pour
nous le caractère du vrai en toutes choses, \struck{chose} nous
\struck{pouvons} \uncertain{pourrons}

\page{10}
\note{en tête, à gauche, le chiffre « 2 » à l'encre.}
\struck{\ill{}} parvenir à reconnoître \struck{comment} que dans les divers
genres d'études, nos recherches, \struck{sont} dirigées vers un même but,
emploient des procédés qui sont aussi les mêmes.

Et, en effet, s'agit-il du plan d'un ouvrage, de l'argument d'un poëme,
l'esprit exige de la clarté ; il veut que les diverses parties soient liées
entre elles avec assez d'art pour que leur rapport s'apperçoive d'un coup
d'\uncertain{œil} ; il demande un ordre facile à saisir ; il se complait dans la
simplicité, source de l'élégance en tout genre. L'emploi du merveilleux est
soumis aux mêmes règles. L'imagination peut adopter d'ingénieuses fictions ;
mais alors un certain module intellectuel remplace ce qui manque à la réalité
des objets. Les oracles du goût et les arrêts de la raison se ressemblent :
l'ordre, la proportion et la simplicité, \struck{remplaçant ce qui manque à la
réalité des objets} ne cessent pas d'être des nécessités intellectuelles. Les
sujets sont différens ; mais le jugement est constamment appuyé par ce type
universel qui appartient également et au beau et au vrai.

Voulons-nous connoître les êtres naturels ? nous les classons suivant nos
convenances ; \& la notion méthodique des genres et des espèces imprime à
l'histoire naturelle le cachet de l'esprit de l'homme.

A l'égard des sciences exactes, le sentiment d'ordre et de proportion, qui
partout ailleurs guide ou le goût ou la raison, fait place à la connoissance
certaine d'un ordre déterminé, de proportions connues et mesurables. On diroit
que, muni d'un instrument nouveau, l'intelligence humaine a renoncé à sa marche
accoutumée. La ressemblance à son modèle intérieur

\note{le mot biffé en tête de page est noirci et illisible. La phrase se poursuit en tête de la vue 13 ; la vue 11 est un feuillet d'addition.}

\page{11}\folio{5r}
\note{feuillet plus petit, monté ; une croix « + » à l'encre dans l'angle supérieur gauche, et le paragraphe ouvert par un crochet. Le texte n'enchaîne ni sur la fin de la vue 10 ni sur le début de la vue 13 : c'est une addition, que la croix renvoie au bas de la vue 13, où la même croix est tracée sous le dernier alinéa.}

Ne nous en étonnons pas ; l'esprit humain obéit à des lois ; elles sont celles
de sa propre existence ; elles lui fournissent une mesure commune entre toutes
les existences qu'il conçoit en dehors de la sienne ; elles deviennent
nécessairement le mobile de tous ses travaux, la source de tous ses plaisirs.
Et, en effet, un trait de génie, un trait d'éloquence, dans les sciences, dans
les beaux arts, dans la \struck{l'éloquence} littérature, nous plaît par une
seule et même raison : il dévoile à nos yeux \struck{une} une foule de rapports
que nous n'avions pas encore apperçus. — nous nous trouvons \struck{conduits}
transportés tout d'un coup dans une région élevée, d'où nous découvrons un
ordre \struck{nouveau} inattendu d'idées ou de sentimens ; le plaisir de la
surprise émeut notre ame ; elle rend un hommage involontaire à son bienfaiteur ;
et \struck{elle trouve encore un plaisir nouveau dans} cet hommage même est
\struck{pour elle la source d'un} encore pour elle un plaisir nouveau.

\note{« transportés » est écrit sous « tout d'un coup », tous deux dans l'interligne au-dessus du mot biffé ; « littérature » et « inattendu » sont également dans l'interligne.}

\page{13}\folio{6r}
\note{en tête, à gauche, un petit « 3 » à l'encre. La première phrase achève celle de la fin de la vue 10.}
n'est plus pour elle le caractère du vrai ; elle l'atteint de plus près ;
l'objet de ses études remplit au plus haut degré les conditions qu'elle cherche
partout ailleurs ; et son attention fixée sur cette heureuse réalisation y est
absorbée toute entière.

Sans doute, l'impression produite par la lecture d'un ouvrage d'imagination ne
ressemble pas à celle qui résulte de l'étude d'un traité de géométrie ; sans
doute aussi, certains esprits admirateurs des riantes images, s'abandonnant
uniquement à ce goût, deviendront tout à fait incapables d'application, tandis
que d'autres esprits exclusivement livrés à la contemplation de la vérité
démontrée demeureront distraits ou incertains lorsqu'ils ne rencontreront pas
une évidence complète.

Ne nous pressons pourtant point de conclure qu'il n'existe aucun lien commun
entre des œuvres qui semblent d'abord si différentes. Assistons à leur création,
et nous reconnoîtrons bientôt que l'esprit humain est guidé dans toutes ses
conceptions par la prévision de certains résultats, vers lesquels se dirigent
tous ses efforts.

En observant la manière dont il procède, nous verrons qu'il agit toujours
suivant une méthode constante ; et, après avoir suivi les différentes époques de
la composition, il deviendra évident que la littérature la plus \ill{}, et les
découvertes \uncertain{d\struck{ont}} \struck{ce genre} s'enrichit
\uncertain{\struck{des}} sciences ont été inspirées par un sentiment d'ordre et
de proportion qui est le régulateur de tout mouvement intellectuel.

\note{passage corrigé en fin de ligne et au début de la suivante : après « découvertes » un mot commençant par « d » est recouvert ; au début de la ligne suivante deux mots biffés (lus « ce genre », avec doute), « s'enrichit » écrit au-dessus, puis un mot surchargé devant « sciences ». La leçon définitive n'est pas établie. Sous le dernier alinéa, une croix « + » : c'est le renvoi du feuillet collé de la vue 11, qui porte la même croix.}

\page{15}\folio{7r}
\note{en tête, à gauche, le chiffre « 5 » à l'encre. Les deux premiers alinéas sont ouverts par un crochet.}

Voyons d'abord quel est le caractère des premiers essais.

Le sujet est choisi ; les idées se présentent en foule à l'imagination du
poëte ; il reste quelque tems incertain ; une multitude de ressorts différens
semblent pouvoir donner la vie à sa composition ; il \struck{\ill{}} en suit le
développement, puis il y renonce ; il fait un choix nouveau ; son mécanisme se
complique ; il n'en est pas content ; il s'arrête, \struck{bientôt} il revient
sur ses pas. Du milieu de cette lutte tumultueuse entre des projets contraires
surgit enfin une idée simple : soit qu'elle ait déjà été entrevue, soit qu'elle
se présente à lui pour la première fois, l'auteur sent que cette idée est celle
qu'il avoit cherchée.

Une remarque, un fait inattendu, donne-t-il lieu à des recherches nouvelles ?
Le géomètre, après avoir mûrement examiné tout ce qui, dans la science déjà
faite, peut lui prêter secours, circonscrit le sujet qu'il va traiter. —
Bientôt il entrevoit des résultats qu'il ne peut encore atteindre ; son
imagination s'élance, pour les saisir, dans les routes qu'elle s'est frayées ;
il craint de s'être égaré, il doute de ses premiers aperçus, il rétrograde et
cherche à ressaisir les indications qui l'avoit d'abord guidé. Un grand nombre
d'idées se sont jointes à celles qui furent les premières, elles compliquent le
sujet, partagent l'attention et suspendent le jugement : mais à travers ce
chaos de pensées diverses, le génie des sciences distingue une idée simple ; son
choix est irrévocablement fixé ; il sait que cette idée sera féconde.

\note{« -t-il » est ajouté dans l'interligne au-dessus de « donne ». « l'avoit d'abord guidé » : ainsi au manuscrit.}

Examinons à présent de quelle manière les travaux commencés vont être exécutés.

\page{16}
\note{en tête, à gauche, le chiffre « 6 » à l'encre. Le premier alinéa, jusqu'à « beauté véritable. », est barré d'une grande croix et en outre biffé ligne à ligne ; il est récrit aussitôt dessous, ouvert par un crochet.}

\struck{En traçant le plan qu'il va suivre, une idée principale guide la pensée
du poëte elle imprimera à son travail cette l'unité d'intérêt et d'action,
source de toute beauté véritable, sans laquelle il ne sauroit satisfaire à ce
besoin d'ordre et de proportion \ill{} que nous avons reconnu se fait sentir dans
tous les genres de composition ; il aura rempli les conditions \ill{} sans
lesquelles il ne pourroit atteindre ni à la clarté \ill{} de rien produire
\ill{} d'atteindre une beauté véritable.}

En traçant le plan qu'il \struck{va} doit suivre, le poëte ne perdra jamais de
vue l'idée principale dont il a fait choix. elle donnera à son travail l'unité
d'intérêt et d'action, source de toute beauté véritable : \struck{et en effet
c'est avec au moyen de cette unité qu'on peut satisfaire au besoin il aura
satisfait à \ill{} se complaira à en suivre la développement car} elle lui
offrira le moyen \struck{étoit faite} \uncertain{au} \struck{ce} besoin d'ordre
et de proportion, \struck{dont le} qu'un sentiment \struck{le procédé la règle
le précepte toute la règle \ill{} dicté la loi} universel a placé au premier
rang \struck{parmi} entre les \struck{\ill{}} préceptes \struck{de la raison et}
du goût et de la raison. Il se complaira à en suivre le développement.

\note{les corrections interlinéaires de ce passage se chevauchent : au-dessus de « en effet » un mot biffé, lu \uncertain{amenaient} ; au-dessus de « a » biffé, « et à ce », biffé ; « partage » dans l'interligne, sans place assignée ; « marche », biffé, au-dessus de « développement ». « lui offrira le moyen » est écrit au-dessus de « étoit faite », biffé ; entre « moyen » et « besoin » la phrase reste incomplète au manuscrit (le verbe manque).}

De son côté le géomètre porte une attention soutenue vers l'idée heureuse qui
dirige ses recherches ; toutes les forces \struck{intellectuelles} de son
intelligence vont être employées à \uncertain{dérouler} la chaîne des vérités
contenues dans cette vérité première ; et sans doute l'unité de composition
\struck{ne peut être nulle} ne sera nulle part ailleurs aussi sensible.
\struck{il n'a} L'ordre de son travail \struck{est} sera déterminé ; il ne
pourroit l'intervertir. \struck{sans nuire à la clarté} L'évidence est pour lui
la condition du succès ; \struck{après} il choisit la méthode qu'il \struck{doit
lui} croit propre à l'y conduire ;

\note{l'alinéa est ouvert par un crochet. Dans la marge de gauche, en regard des deux dernières lignes, deux mots biffés, lus \uncertain{avant} et « et ».}

\page{17}\folio{8r}
\note{en tête, à gauche, le chiffre « 7 » à l'encre. La première phrase achève celle de la fin de la vue 16.}
\struck{l'autorité de la démonstration succèdera bientôt il peut se flatter devoir
bientôt l'autorité de sa démonstration succèdera aux simples prévisions de son
jugement et il} Et il entre alors avec joie dans la carrière \struck{qui
s'ouvre} ouverte à ses espérances.

\note{« ou bientôt », biffé, dans l'interligne au-dessus de « succèdera » ; « Et il » au-dessus de « et il » biffé ; « alors » au-dessus de « avec », et sous ce mot « ensuite » entouré d'un trait ; « ouverte » au-dessus de « s'ouvre » biffé.}

\note{dans la marge de gauche, devant l'alinéa suivant, un nom d'une autre écriture, lu \uncertain{Cherenoz}, comme « Gutmann » à la vue 9. L'alinéa est ouvert par un crochet.}

Les auteurs dont nous comparons les travaux, ont franchi les premières
difficultés ; ils \struck{observeront} observé entre les divisions qu'ils
\struck{ont tracés} viennent d'adopter, cette juste proportion d'étendue
respective qui, sans nuire au sentiment de la continuité, offre à l'attention
les repos dont elle a besoin.

Pour remplir ensuite les cadres qu'ils ont tracés, ils s'abandonneront encore
une fois aux inspirations de leur génie. mais, à présent que les limites du
sujet sont parfaitement déterminées, ils n'auront plus à craindre de s'égarer,
l'un dans le champ immense d'une imagination fertile en inventions, l'autre
dans \struck{cet} cet océan des possibilités d'où on aborde avec tant de
difficulté sur le terrain ferme de la vérité démontrée.

\struck{Comment l'époque ou et époque des travaux} \uncertain{Il se présentera
toujours} encore des idées \struck{dans le cours du travail} qui, bien que nées
du sujet, nuiroient cependant ou à la rapidité ou à la clarté de son
développement. S'ils mettoient trop de soin à éviter une telle surabondance
d'invention nos auteurs craindroient d'arrêter l'élan de leurs pensées. plus
tard ils reverront leurs premières ébauches, et n'y conserveront plus que les
\struck{seuls} traits nécessaires. Changeant alors de rôle, ils deviennent les
juges de leurs propres ouvrages.

\note{« tant » est dans la marge, « de » dans l'interligne. Au-dessus de la ligne biffée « Comment l'époque… » deux groupes de mots biffés, illisibles ; « dans le cours du travail » est écrit sous « encore des idées » et entouré d'un trait, sans place sûre dans la phrase ; « encore des idées » est souligné. Le début de la phrase (« Il se présentera toujours ») est une lecture incertaine d'une ligne surchargée.}

Ils examinent d'abord la marche des idées ; celles qui pourroient,

\note{l'alinéa est ouvert par un crochet ; la phrase se poursuit à la vue 18.}

\page{18}
\note{en tête, à gauche, un chiffre à l'encre presque entièrement recouvert d'une tache, sans doute « 8 » dans la suite des feuillets voisins. La première ligne continue la phrase de la fin de la vue 17.}
d'un côté, partager l'intérêt, de l'autre, suspendre l'attention et détruire
ainsi l'unité de composition, seront écartées du lieu où elles se trouvent.
Elles iront enrichir soit de gracieux épisodes, soit de savantes annotations ;
ou, si, trop éloignées du sujet qui les a fortuitement amenées, elles ne
peuvent être convenablement placées dans l'ouvrage même, elles deviendront
peut-être l'origine d'une production nouvelle. Ainsi la branche développée
dans la saison actuelle offre quelquefois le rudiment d'une végétation
prochaine.

Les différentes parties du \uncertain{style} seront ensuite l'objet d'un autre
genre de corrections. L'homme de lettres s'occupera du choix des mots, de leur
arrangement, de l'harmonie du vers, ou de celle de la phrase. Un grand nombre de
convenances difficiles à concilier, seront soumises au jugement du goût, du goût
tantôt si prompt à décider, tantôt si lent à prononcer, dont les opérations
échappent souvent à l'attention, mais qui pourtant agit toujours conformément
aux règles de la raison lors même qu'il semble ne reconnoître d'autres lois que
ses propres caprices.

\note{l'alinéa est ouvert par un crochet ; le mot « style » est en partie couvert par une tache d'encre.}

La langue des calculs peut donner lieu à des \struck{observations} corrections
\struck{semblables} qui lui sont propres ; car elle a aussi son style, et tous
les auteurs ne l'écrivent pas avec le même degré de perfection.

\note{« corrections » et « qui lui sont propres ; car » sont écrits dans l'interligne ; dans ce dernier ajout, un mot surchargé devant « propres », lu « sont ».}

\page{19}\folio{9r}
\note{en tête, à gauche, le chiffre « 9 » à l'encre ; au centre, d'une autre écriture, un nom lu \uncertain{Deraihe}, comme « Gutmann » à la vue 9 et « Cherenoz » à la vue 17.}

Au choix des mots correspond celui des caractères. \struck{La vérité} Ceux-ci
sont tellement conventionnels, qu'il faut, dans chaque occasion,
\struck{avertir} exprimer quelle valeur on leur attribue : cependant leur emploi
est assujetti à certaines convenances qui ne tiennent pas uniquement aux
habitudes consacrées. Les formules remplacent la phrase ; elles peuvent être
plus ou moins élégantes. L'analyse parle aux yeux ; ainsi, au lieu de l'harmonie
ou de l'accord entre les sons, elle \struck{est assujettie à} doit présenter
entre ses divers élémens des rapports d'ordre et de simplicité faciles à saisir
au premier coup d'œil. \struck{il en résulte pour} Les personnes
\struck{exercées} initiées à ce genre de discours \struck{trouvent dans l'analyse
les \ill{} une sorte d'attrait \ill{}} trouvent bien certainement dans la
contemplation des formules une sorte \struck{de} de charme, qui les entraîne
vers l'étude. \struck{des ouvrages q} Et si les bons auteurs sont doués d'une
finesse de tact, qui \struck{leur avertit des formes} \uncertain{leur} fait
choisir, \struck{les formes et} entre ces formules, celles qu'il convient
d'écrire, tandis que d'autres seront seulement indiquées ; \struck{à cet égard
aussi tantôt les raisons leur décision est quelquefois instantanée} si leurs
décisions sont \struck{aussi} tantôt rapides, tantôt lentes \struck{comme} et
réfléchies ; c'est que le tact dont nous parlons n'est en effet autre chose que
le goût \struck{appliqué à des matières} appliqué à des \struck{matières} objets
qu'on semble avoir crus étrangers à son empire.

\note{les corrections des lignes « Les personnes… de charme » se superposent en trois couches dans l'interligne et à droite, reliées par des traits ; la leçon retenue est celle qui n'est pas biffée. « rapports » et « ce charme » sont écrits dans la marge de gauche, le second pour remplacer « de charme » ; « Et si » est écrit sous « des ouvrages » biffé. Le mot devant « fait choisir » est couvert d'une tache.}

Nous venons de voir combien les productions intellectuelles les plus diverses
ont entr'elles de ressemblances véritables ; \struck{et} comment un sentiment
d'ordre et de proportion, après avoir présidé aux inspirations du génie guide
leur emploi et se fait encore sentir dans les dernières corrections de
l'ouvrage achevé ; \struck{tellement que dans des genres de perceptions}

\note{la dernière ligne, biffée, s'interrompt au bas du feuillet ; la vue 20 commence un nouvel alinéa.}

\page{20}
\note{en tête, à gauche, le nombre « 10 » à l'encre ; aucun numéro de feuillet n'est visible dans l'angle droit. Le haut de la page est très raturé ; « mais » est écrit dans la marge devant « si ».}

Mais si la \struck{opérations} marche de l'esprit \struck{sont} est partout la
même, \struck{leurs} les objets \struck{dont il s'occupe} \struck{peuvent être
tellement différens qu'ils écartent, au premier aperçu, toute idée de
rapprochement entre} qu'il \struck{envisage} peut envisager, sont d'une variété
\struck{à} infinie. \struck{leur différences frappent au premier coup d'œil
beaucoup plus} Au premier coup d'œil, \struck{leurs différences sont} ce qui
tient à \struck{leur nature} \ill{} cette variété dont \add{nous avons}
parlé frappe \struck{bien plus} \struck{\ill{} davantage} plus que l'identité des
rapports \struck{que nous avons reconnus}.

\note{la phrase « Au premier coup d'œil… » est reconstituée à partir de corrections superposées dans les deux interlignes : au-dessus, « ce qui tient à », « leur nature » biffé, un mot surchargé (\ill{}), « cette variété » souligné, puis « bien plus » biffé ; au-dessous, « dont » et « parlé ». L'ordre retenu est une conjecture ; « nous avons » n'est pas écrit : c'est un ajout de l'éditeur.}

Aussi des opérations intellectuelles qui \struck{sont semblables en tout}
\uncertain{confondues}, les mêmes, ont-elles reçu \struck{des} divers noms,
\struck{différens} suivant \struck{qu'elles s'appliquent à} la nature des sujets
auxquels elles s'appliquent. \struck{Cette} La différence dans les mots,
différence d'autant plus naturelle que chacune des branches de nos
connoissances \struck{ont} a été pendant longtems, pour ainsi dire, exclusive
\struck{l'une de l'autre} de toutes les autres, tend à perpétuer l'opinion d'une
séparation réelle entre les facultés \struck{intellectuelles ou} \uncertain{de
l'esprit} comme si, par exemple, l'allégorie elle même n'étoit pas assujettie aux
préceptes de la raison, et si la découverte d'une loi de la nature avoit pu se
passer du secours de l'imagination.

\note{« confondues » et « les mêmes » sont dans l'interligne au-dessus de « sont » et « en tout », biffés ; au-dessus de « intellectuelles », un mot biffé, et au-dessous « de l'esprit ».}

Sans doute le poëte \struck{ne} ne nous rendra pas compte des discussions
pleines de finesse \struck{auxquelles il a soumis} qui ont précédé
\struck{le choix} \uncertain{l'adoption} des emblêmes qu'il \struck{a employés
emploie} \uncertain{a} choisis : l'homme de génie qui a surpris un des secrets
de l'ordre naturel ne nous dira pas non plus combien \struck{son imagi} depuis
son imagination s'est égarée \struck{dans le voisinage} autour de la route qui
devoit le conduire à la

\note{au-dessus de « le choix », biffé, un mot lu \uncertain{l'adoption} ; au-dessus de « a employés », « adoptés », biffé, et au-dessous « choisis ». « non plus » est écrit dans la marge de gauche. La phrase s'interrompt au bas de la page et se poursuit au-delà du lot.}

\end{document}
